\hypertarget{prerequisites}{%
\section{Prerequisites}}

Before starting this workshop, ensure you have the following tools
installed and configured.

\hypertarget{required-accounts-permissions}{%
\subsection{Required Accounts \&
Permissions}}

\hypertarget{aws-account}{%
\subsubsection{AWS Account}}

\begin{itemize}
\tightlist
\item
  \textbf{AWS Account} with administrative access or permissions for:

  \begin{itemize}
  \tightlist
  \item
    VPC, EC2, Subnets, Route Tables, Internet Gateway
  \item
    RDS (Aurora PostgreSQL Serverless v2)
  \item
    ECS (Fargate, Clusters, Task Definitions, Services)
  \item
    ECR (Repositories)
  \item
    Lambda (Functions, Layers, Permissions)
  \item
    SQS (Queues)
  \item
    EventBridge (Rules, Schedules, Targets)
  \item
    API Gateway (REST APIs)
  \item
    Bedrock (Model access: \texttt{amazon.titan-embed-text-v2:0})
  \item
    IAM (Roles, Policies)
  \item
    CloudWatch (Log Groups, Metrics)
  \item
    CloudFormation (Stacks)
  \end{itemize}
\end{itemize}

\begin{quote}
\textbf{Tip:} Use an IAM user with \texttt{AdministratorAccess} for the
workshop, or ensure your role has the above permissions.
\end{quote}

\hypertarget{external-api-keys-for-rag}{%
\subsubsection{External API Keys (for
RAG)}}

\begin{itemize}
\tightlist
\item
  \textbf{Groq API Key} --- Get from
  \href{https://console.groq.com/}{console.groq.com} (Free tier
  available)
\item
  \textbf{Google Gemini API Key} --- Get from
  \href{https://aistudio.google.com/}{aistudio.google.com} (Free tier
  available)
\end{itemize}

\hypertarget{required-tools}{%
\subsection{Required Tools}}

\hypertarget{aws-cli-v2}{%
\subsubsection{1. AWS CLI v2}}

\emph{(Mã nguồn chi tiết đã được lược bỏ để báo cáo ngắn gọn. Vui lòng
xem trong source code dự án)}

\textbf{Configure AWS CLI:}

\begin{Shaded}
\begin{Highlighting}[]
\ExtensionTok{aws}\NormalTok{ configure}
\CommentTok{\# Enter: AWS Access Key ID, Secret Access Key, Region (e.g., ap{-}southeast{-}2), Output format (json)}
\end{Highlighting}
\end{Shaded}

\hypertarget{terraform-1.5.0}{%
\subsubsection{2. Terraform \textgreater= 1.5.0}}

\emph{(Mã nguồn chi tiết đã được lược bỏ để báo cáo ngắn gọn. Vui lòng
xem trong source code dự án)}

\hypertarget{docker-docker-compose}{%
\subsubsection{3. Docker \& Docker
Compose}}

\emph{(Mã nguồn chi tiết đã được lược bỏ để báo cáo ngắn gọn. Vui lòng
xem trong source code dự án)}

\hypertarget{python-3.10}{%
\subsubsection{4. Python 3.10+}}

\emph{(Mã nguồn chi tiết đã được lược bỏ để báo cáo ngắn gọn. Vui lòng
xem trong source code dự án)}

\hypertarget{git}{%
\subsubsection{5. Git}}

\emph{(Mã nguồn chi tiết đã được lược bỏ để báo cáo ngắn gọn. Vui lòng
xem trong source code dự án)}

\hypertarget{code-editor-recommended-vs-code}{%
\subsubsection{6. Code Editor (Recommended: VS
Code)}}

\emph{(Mã nguồn chi tiết đã được lược bỏ để báo cáo ngắn gọn. Vui lòng
xem trong source code dự án)}

\hypertarget{project-setup}{%
\subsection{Project Setup}}

\hypertarget{clone-the-repository}{%
\subsubsection{Clone the Repository}}

\begin{Shaded}
\begin{Highlighting}[]
\FunctionTok{git}\NormalTok{ clone }\OperatorTok{\textless{}}\NormalTok{your{-}repo{-}url}\OperatorTok{\textgreater{}}\NormalTok{ AWS{-}Projects}
\BuiltInTok{cd}\NormalTok{ AWS{-}Projects}
\end{Highlighting}
\end{Shaded}

\hypertarget{python-virtual-environment}{%
\subsubsection{Python Virtual
Environment}}

\begin{Shaded}
\begin{Highlighting}[]
\ExtensionTok{python3} \AttributeTok{{-}m}\NormalTok{ venv venv}
\BuiltInTok{source}\NormalTok{ venv/bin/activate  }\CommentTok{\# Linux/macOS}
\CommentTok{\# venv\textbackslash{}Scripts\textbackslash{}activate   \# Windows}

\ExtensionTok{pip}\NormalTok{ install }\AttributeTok{{-}{-}upgrade}\NormalTok{ pip}
\ExtensionTok{pip}\NormalTok{ install }\AttributeTok{{-}r}\NormalTok{ requirements.txt}
\end{Highlighting}
\end{Shaded}

\hypertarget{environment-variables}{%
\subsubsection{Environment Variables}}

Copy the example file and fill in your values:

\begin{Shaded}
\begin{Highlighting}[]
\FunctionTok{cp}\NormalTok{ .env.example .env}
\end{Highlighting}
\end{Shaded}

Edit \texttt{.env} with your credentials:

\emph{(Mã nguồn chi tiết đã được lược bỏ để báo cáo ngắn gọn. Vui lòng
xem trong source code dự án)}

\hypertarget{terraform-variables}{%
\subsubsection{Terraform Variables}}

Create \texttt{terraform.tfvars}:

\begin{Shaded}
\begin{Highlighting}[]
\FunctionTok{cat} \OperatorTok{\textgreater{}}\NormalTok{ terraform.tfvars }\OperatorTok{\textless{}\textless{} EOF}
\StringTok{db\_password     = "your\_secure\_password\_123!"}
\StringTok{db\_user         = "postgres"}
\StringTok{qdrant\_api\_key  = ""  \# Leave empty for AWS deployment}
\StringTok{qdrant\_host     = ""  \# Leave empty for AWS deployment}
\StringTok{model\_1\_api\_key = "gsk\_xxxxxxxxxxxxxxxxxxxx"  \# Groq}
\StringTok{model\_2\_api\_key = "gsk\_xxxxxxxxxxxxxxxxxxxx"  \# Groq backup}
\StringTok{model\_3\_api\_key = "AIza\_xxxxxxxxxxxxxxxxxxxx" \# Gemini}
\OperatorTok{EOF}
\end{Highlighting}
\end{Shaded}

\begin{quote}
\textbf{Security:} Never commit \texttt{.env} or
\texttt{terraform.tfvars} to git. They are in \texttt{.gitignore}.
\end{quote}

\hypertarget{verify-setup}{%
\subsection{Verify Setup}}

Run the verification script:

\begin{Shaded}
\begin{Highlighting}[]
\FunctionTok{make}\NormalTok{ verify}
\end{Highlighting}
\end{Shaded}

Or manually check each tool:

\begin{Shaded}
\begin{Highlighting}[]
\ExtensionTok{aws}\NormalTok{ sts get{-}caller{-}identity  }\CommentTok{\# Should show your AWS account}
\ExtensionTok{terraform}\NormalTok{ version            }\CommentTok{\# Should be \textgreater{}= 1.5.0}
\ExtensionTok{docker}\NormalTok{ run hello{-}world       }\CommentTok{\# Should print "Hello from Docker!"}
\ExtensionTok{python3} \AttributeTok{{-}c} \StringTok{"import scrapy; print(\textquotesingle{}Scrapy OK\textquotesingle{})"}
\end{Highlighting}
\end{Shaded}

\hypertarget{aws-region}{%
\subsection{AWS Region}}

This workshop uses \textbf{ap-southeast-2 (Sydney)} by default. To
change: 1. Update
\texttt{provider\ "aws"\ \{\ region\ =\ "your-region"\ \}} in
\texttt{main.tf} 2. Update \texttt{AWS\_REGION} in \texttt{.env} 3.
Ensure Bedrock model \texttt{amazon.titan-embed-text-v2:0} is available
in your region (check
\href{https://docs.aws.amazon.com/bedrock/latest/userguide/models-regions.html}{Bedrock
regions})

\hypertarget{bedrock-model-access}{%
\subsection{Bedrock Model Access}}

Enable model access in AWS Console: 1. Go to \textbf{Amazon Bedrock}
\(\rightarrow\) \textbf{Model access} 2. Click \textbf{Manage model
access} 3. Enable \textbf{Amazon Titan Embeddings G1 - Text v2}
(\texttt{amazon.titan-embed-text-v2:0}) 4. Wait for status to show
``Access granted''

\begin{center}\rule{0.5\linewidth}{0.5pt}\end{center}

\textbf{Next:} \href{5.3-Infrastructure/}{Infrastructure as Code
(Terraform)}
